\documentclass[11pt]{article}
    \usepackage[utf8]{inputenc}
    \usepackage[T1]{fontenc}
    \usepackage[margin=1in]{geometry}
    \usepackage{hyperref}
    \usepackage{enumitem}
    \title{Block Coding with Bees (Intro. Block Coding & Loops)}
    \author{Piter Garcia}
    \date{March 20, 2026}
    \begin{document}
    \maketitle
    \section*{Quick Scan}
    \begin{itemize}[leftmargin=*]
    \item Grade: Grade 2
    \item Workbook sessions: Session 22
    \item Class blocks: Day 1 10:45-11:35 Callon, Day 1 2:15-3:05 Henchen, Day 4 10:45-11:35 Barthelman
    \item Reference file: block-coding-with-bees-intro-block-coding-and-loops.bib
    \end{itemize}
    \section*{School Planning Goals and Inclusion}
    \begin{itemize}[leftmargin=*]
    \item What is expected: I can decompose a problem into smaller steps, in multiple ways | | I can identify the essential information needed to get the Bee to the hive. | | I can code using fixed and changing pieces of information. | | I can document my plan to get the bee to the hive using block-coding
    \item What we achieved: No direct transcript match is attached yet. Treat this lesson as workbook-backed and enrich it when a matching classroom transcript is available.
    \item What needs improvement: stronger proof, cleaner assessment notes, and tighter lesson artifacts.
    \item How to improve: add transcript proof, one saved artifact, and clearer success criteria.
    \item Inclusion focus: use visual modeling, chunked directions, predictable routines, and flexible participation roles.
    \end{itemize}
    \section*{Official References}
    \begin{itemize}[leftmargin=*]
    \item \url{https://csteachers.org/k12standards/interactive/}
\item \url{https://code.org/curriculum/elementary-school}

    \end{itemize}
    \end{document}
